\subsection{Pipeline Generality Across Domains}
\label{sec:cross_domain}

\noindent\textbf{Is the pipeline bearing-specific?}
To demonstrate that \neuroplc's compiler and verification
pipeline is domain-agnostic---depending only on the KAN
architecture, not on data provenance---we apply the
identical pipeline to a fundamentally different task:
MNIST handwritten digit classification (digits 0--3).
Raw $28\times28$ pixel images are reduced to 28-D
features via PCA, then z-score normalized and clipped
to $[-3,3]$---matching the preprocessing contract of the
CWRU bearing pipeline. The same KAN $[28,16,4]$
architecture is trained, compiled, and verified with
\textbf{zero modification} to the compiler or verification
toolchain.

\begin{table}[ht]
\centering
\caption{Cross-Domain Pipeline Generality: Identical compiler
and verification pipeline applied to image classification
(MNIST) and vibration-based fault diagnosis (CWRU).}
\label{tab:cross_domain}
\begin{tabular}{@{}lcc@{}}
\toprule
\textbf{Metric} & \textbf{MNIST (Image)} & \textbf{CWRU (Vibration)} \\
\midrule
  Domain & Handwritten digits & Bearing faults \\
  Features & PCA(784$\to$28) & RCMDE $+$ RCHFDE \\
  Test accuracy & $0.9859$ & $0.9875$ \\
\midrule
  Per-function verified & $512/512$ & $512/512$ \\
  DA bound $\Delta_{\text{DA}}$ & $0.1199$ & $0.1196$ \\
  DA/IA tightening & $4.8\times$ & $2.1\times$ \\
  Certificate valid & Yes & Yes \\
\bottomrule
\end{tabular}
\end{table}

\textbf{Key Insight.} The compiler and verification
pipeline require only that the model be a KAN---the
data domain is irrelevant. The B-spline functions are
verified against their LUT approximations using the same
$M_2 h^2/8$ bound regardless of whether the input features
represent pixel intensities or vibration amplitudes.
The DA composition rules depend only on the weight matrix
structure, not on data semantics. This domain independence
is a direct consequence of the SVNN framework
({\S}\ref{sec:svnn}): the three sufficient conditions are
architectural properties of KAN, not properties of the
training data or task.
